\section{Use the chapter package for feedback and submission}
The chapter package is \texttt{034-lossless-video}. It contains this native reading,
all nine construction stages, numeric input generators, public completed-interface
checks and separate grading storage. Guided work follows the 122 steps; free-writing
work can use another private layout or encoding policy while preserving the named
interfaces and their state contracts. The package contains no completed assembly
codec or quiz answer keys.

Keep the two early working threads separate. Begin XOR work with
\path{construction/video-construction/kernel.asm}. Begin record decoding with
\path{construction/video-record-construction/main.asm}, then attach indexing,
loading, seeking and full-record encoding in their chapter order. The delta stage
imports your XOR working file with its \texttt{--kernel-source} argument. File
encoding and playback then extend that combined source. Each attachment retains
your earlier instructions; keep those working files rather than replacing completed
regions with a fresh template. Use a stage's preparation helper with the working
source for that stage. Some earlier example labels are namespaced in the combined
source, so their original preparation helpers are not interchangeable with later
ones.

The optional \texttt{starter/main.asm} contains all unfinished regions together.
For a smaller free-writing file, \texttt{starter/interfaces.asm} supplies only the
public entries and an unfinished main. Completed checks call those entries directly;
use the full caller when you want the interactive file round trip.
The default eight-frame example in \texttt{starter/main.asm} uses smaller provided arenas: 262\,144 bytes each
for staged file storage and packed-source allocation. Those sizes fit the supplied
input and avoid reserving two maximum arrays during ordinary development. They do
not reduce a function's input limits: the grader passes its own argument windows
when checking maximum-size work. The functions must honor those arguments rather
than relying on example-buffer addresses or a fixed eight-frame count.

After completing the construction sequence, use
\texttt{python3 construction/compact.py --source working.asm --output submit.asm}
to produce the same small default profile while retaining your instructions and
other data. The helper reduces only unused tails of the two supplied arenas and
updates their encoder and loader capacity controls. It refuses changed markers,
an input too large for that profile, and an existing output file. Inspect the result
before submission. A custom free-writing project can keep its own storage layout;
it need not use this helper.

Run a completed public domain with \texttt{python3 grade.py submit.asm --domain
record-encode --cli /path/to/cli.jar --java /path/to/java}. Omit \texttt{--domain}
to run all eleven domains: XOR, record application, indexing, loading, seeking,
full-record encoding, record selection, whole-file encoding, storage, playback and
XOR reuse. These finite checks use published numeric fixtures and deliberately bad
inputs. The runner adds \texttt{checks/harness.asm} as another assembly source.
That file supplies buffers in the separate wide-data region and a 32\,768-byte
private stack. It supplies no algorithm. Reserve the \texttt{videoCheck\_} and
\texttt{videoPrivate\_} names for grading, and keep your own static data in the
ordinary region. Constants and different private frames are permitted.

The instructor suite adds larger and private cases. Its largest file/source checks
use ordinary-data grading buffers; keep student-owned ordinary data within the first two MiB for
that profile. The compact supplied example is below that budget. The supplied
private stack and caller-owned work buffers provide the storage needed by every
interface. The reserved boundary arena begins at \texttt{0x10200000}. This is a finite simulator allocation budget, not a requirement for a
particular stack frame, scratch-register choice or label layout. Submit your actual
implementation and its required constants; do not concatenate all historical
example files or include the grading harness in your own source.

An unfinished routine is expected to fail completed checks. During guided work,
stop at the step's named checkpoint and compare its observation; use the completed
domain after filling the stage. These are different questions. A checkpoint checks
illustrated intermediate state, while completed feedback checks the callable result.
Instructor controls can select quiz requirements, hint availability, skipping and
mode changes. The package records those controls; automatic remote delivery and
progression enforcement remain future work in the separate authoring plugin.

Retain evidence from both successful and failed calls. Record the entry, input,
selected byte order, status/count and the pixels or fields that should remain.
For a decoder failure, keep an earlier accepted image through the rejected retry.
For a file-encoding failure, distinguish retained final context from temporary
staging and working images, which may change. A failed truncate-write can leave an
empty target or transferred prefix. A playback clock failure can leave accepted
pictures visible and returns their actual count. These are different publication
contracts; a single rule that clears every output on every failure is incorrect.

Equivalent full/delta choices, additional permitted full frames and valid alternative
command boundaries are accepted by completed encoding checks. Required periodic
full frames, exact reconstruction, CRC coverage and complete file boundaries remain
mandatory. A smaller record buffer may succeed with a compact choice or reject a
selected representation that does not fit. Instruction counts are measurements with
named boundaries, rather than required instruction sequences for a correct solution.

Playback feedback uses scripted time and nonblocking sleep. It inspects every accepted
picture at public syscall boundaries, including unsigned clock transitions, early
wakeups, lateness and failure after visible progress. The ordinary program instead
uses the real host clock and sleep service. Its lateness may differ between runs,
and no memory-store assertion establishes an IDE repaint time. Report deterministic
arithmetic and actual host observations separately.

The package includes planned IntelliJ Light captures for XOR lanes, accepted record
state, directory publication, equivalent seek starts, record selection and playback.
Those plans are not completed screenshots. Actual images, native annotations and
coordinate receipts belong to the separate authoring plugin. The initial text and
package draft can be reviewed now; actual captures and independent accessibility
verification remain pending. Substantive editing is the separate later goal.
